\documentclass[a4paper,10pt]{article}

\usepackage[ngerman]{babel}
%\usepackage[top=3cm,bottom=3cm,left=3cm,right=3cm]{geometry} %NICHT IN KEMIE2


\usepackage{microtype} %schöneres typesetting  mit [stretch=10] für nur 1% anstelle von 2%
\usepackage{lmodern} %Latin Modern FONT
\usepackage{pifont} %schönere Font

\usepackage{chemfig}
\usepackage[version=4]{mhchem}
\usepackage{amsmath}
\usepackage{nicefrac}
\usepackage{amssymb}
%\usepackage{mathtools}

\usepackage{titlesec} %Um Section, etc bearbeiten zu können
\usepackage{fancyhdr}
\usepackage{setspace}
\usepackage{enumitem} %enumerate
\usepackage{arydshln} %für dashline
\usepackage{multicol}
\usepackage{lipsum}
\usepackage[labelfont=bf]{caption}

\usepackage[many]{tcolorbox}

%\usepackage{pgfornament} %Scheene Ornamente
\usepackage{witharrows} %Arrows in align-Umgebung
%\usepackage{awesomebox} %Notizboxen, etc
%\usepackage{textpos} %Für zum Beispiel vertikale Texte am Rand


%\usepackage{wrapfig} %Text um Bilder
%\usepackage{varwidth} % Alternative für engere minipage


\usepackage{hyperref}


\tcbuselibrary{theorems}

%\titlespacing*{\section}{0pt}{10mm}{3mm}
%\titlespacing*{\subsection}{0pt}{6mm}{3mm}
%\titlespacing*{\subsubsection}{0pt}{8mm}{3mm}


\renewcommand{\ttdefault}{lmtt} %Schriftart wird geändert zu Latin Modern Typewriter


\hypersetup{hidelinks,
			colorlinks=true,
			linkcolor=black,
			citecolor=miudiutex,
			urlcolor=miugray2
			}

\setlength\fboxsep{0pt}				
\setlength\columnsep{20pt}
\setlength{\parindent}{0pt}
\setcounter{section}{0}






%\definecolor{miudiutex}{HTML}{2f3d39}
\definecolor{miudiutex}{HTML}{1d2926}
\definecolor{miugray}{HTML}{b4cccf}
%\definecolor{miugray2}{HTML}{4d7365}
\definecolor{miugray2}{HTML}{6F8C81}
\definecolor{beige}{HTML}{F4F8D3}
\definecolor{white2}{HTML}{f5f7f8}
\definecolor{red}{HTML}{cf1f5c}
\definecolor{red2}{HTML}{9C191B}
\definecolor{green}{HTML}{00A36C}
\definecolor{delphinidin2}{HTML}{315794}
\definecolor{delphinidin}{HTML}{8B9EBA}
\definecolor{uniblau}{HTML}{004e9f}
\definecolor{unigelb}{HTML}{fcba00}
\definecolor{unigrau}{HTML}{909085}
\definecolor{pelargonidin}{HTML}{b33807}
\definecolor{pelargonidin2}{HTML}{FFDC96}
\definecolor{cyanidin}{HTML}{ba0746}
\definecolor{paeonidin}{HTML}{d15ca6}
\definecolor{petunidin}{HTML}{B88495}
\definecolor{petunitext}{HTML}{963354}
\definecolor{malvidin}{HTML}{8a2d8c}
\definecolor{mytheorembg}{HTML}{F2F2F9}
\definecolor{mytheoremfr}{HTML}{00007B}
\definecolor{uniblau2}{HTML}{27548A}
\definecolor{unigelb2}{HTML}{DDA853}
\definecolor{unidunkelblau2}{HTML}{183B4E}
\definecolor{unibeige}{HTML}{F5EEDC}


\tikzset{>=Latex}

\raggedcolumns %wichtig für Multicolumns
\color{miudiutex}



\usetikzlibrary{chains,
                positioning,
                shapes.multipart,
                }


\setchemfig
{%	
	bond offset=2pt,
	atom sep=15pt, 
	cram width=3pt,
	cram dash width=0.8pt,
	cram dash sep=1.2pt,
	double bond sep=2.5pt, 
	bond style={line width=1.1pt,cap=round},
	baseline=(current bounding box.center),
	bond join=true
}
\setcharge{extra sep=3pt}





%================================
% SYMBOLS
%================================

\newcommand{\RA}{$\rightarrow$~}
\newcommand{\RL}{$\rightleftharpoons$~}
\newcommand{\HARP}{\ding{226}~}
%$\xrightarrow{2}$ %Pfeil mit 2 drüber


%================================
% SHORT COMMANDS (BOXES ; ETC)
%================================


\newcommand{\info}[1]{\begin{note}{\color{miudiutex}#1}\end{note}}
\newcommand{\qsbig}[1]{\begin{qstion}{}{}\begin{center}
			#1 \end{center}\end{qstion}}
\newcommand{\notiz}[1]{\begin{notizz}{}{}\begin{center}
			{\color{miugray2!50!miudiutex}#1} \end{center}\end{notizz}}
\newcommand{\klausur}[1]{\begin{klausurinfo}{}{}\begin{center}
			#1 \end{center}\end{klausurinfo}}
\newcommand{\dfn}[1]{\begin{defin}{}{}\begin{center}
			#1 \end{center}\end{defin}}
\newcommand{\gesetz}[2]{\begin{gestex}{}{}\begin{center}
			#1 \end{center}\end{gestex}}
\newcommand{\gergesetz}[1]{\begin{gesgex}{}{}\begin{center}
			#1 \end{center}\end{gesgex}}
\newcommand{\klaausur}[1]{\begin{klausurinfo}[width=.77\linewidth]{}{}\begin{center}
			\footnotesize #1 \end{center}\end{klausurinfo}}
\newcommand{\winfo}[1]{\begin{wichtig}{}{} \begin{center}
#1 \end{center} \vspace*{0mm} \end{wichtig}}
\newcommand{\zit}[2]{\begin{zitat}{#1}{}\small #2 \end{zitat}}



\newcommand{\unten}{\end{center} \tcblower \begin{center}}


\newcommand*\circled[1]{\tikz[baseline=(char.base)]{
            \node[shape=circle,draw,inner sep=0.5pt,miudiutex] (char) {\tiny #1};}}
            
\newcommand{\miusquare}{{\color{miugray2!50}\scriptsize $\blacksquare$}~}
\newcommand{\miusquarp}{{\color{petunidin!50}\scriptsize $\blacktriangleright$}~}
\newcommand{\niu}{\item[\miusquare]}
\newcommand{\nip}{\item[\miusquarp]}
\newcommand{\mrk}[1]{{\color{petunitext!55!red}\textsc{#1}}}

%%%%% ABSAETZE

\newcommand{\pps}{\par \vspace*{1mm}}
\newcommand{\pp}{\par \vspace*{3mm}}
\newcommand{\ppbig}{\par \vspace*{8mm}}
\newcommand{\pphuge}{\par \vspace*{10mm}}



%================================
% HEADER & FOOTER FUER UNI
%================================
\newcommand{\miusfoot}{
\fancypagestyle{plain}{%
  \fancyhf{}%
  \renewcommand{\footrulewidth}{0.0mm}%
  \fancyfoot[L]{\hspace*{-3cm} {\color{miugray2}\textbf{\thepage}}}%
  \renewcommand{\headrulewidth}{0mm}%
} 
\pagestyle{plain}
}

\newcommand{\miushead}[2]{
\miusquare #1 - 4.Semester - ELW \hfill 
\includegraphics[width=60pt]{{F://LaTeX-Files/shiet/unibo}}
\par \vspace*{-3mm}
\hrulefill
\par \vspace*{1mm}
\par \vspace*{-4mm}
\normalsize	
	\begin{flushright}
	\textbf{\miusquare #2}\par \vspace*{1.5mm}
	\small	
	\textit{Letzte Überarbeitung:} \par 
	\footnotesize \today
	\end{flushright}
	\normalsize
\par \vspace*{-8mm}
}

\newcommand{\sidebar}{
\AddToHook{shipout/background}{
  \begin{tikzpicture}[remember picture,overlay, shift={(current page.south west)}]
   \path[fill=miugray2!15](0,0) rectangle (3.5cm,\paperheight);
%   \path[fill=miugray2!8](0,0) rectangle (3.5cm,\paperheight);
   \path[draw,miudiutex!50] (3.5,0) to (3.5,\paperheight);
%   \path[draw,miudiutex!50,dashed] (3.5,0) to (3.5,\paperheight);  
%   \node[opacity=0.8](a)at(12,10){\includegraphics[width=9cm]{F://LaTeX-Files/pngs/girl01}};
  \end{tikzpicture} 
}
}


\newcommand{\oddsidebar}{
\AddToHook{shipout/background}{
   \put(0in,-\paperheight){%
      \ifodd\value{page}\relax
           \begin{tikzpicture}[remember picture,overlay, shift={(current page.south west)}]
   \path[fill=miugray2!15](0,0) rectangle (3.5cm,\paperheight);
      \path[draw,miudiutex!50] (3.5,0) to (3.5,\paperheight);
  \end{tikzpicture}
      \else
           \begin{tikzpicture}[remember picture,overlay, shift={(current page.south east)}]
   \path[fill=miugray2!15](-3,0) rectangle (3cm,\paperheight);
%      \path[draw,miudiutex!50] (-3.5,0) to (-3.5,\paperheight);
  \end{tikzpicture}
      \fi
   }
}
}
%================================
% ENVIRONMENTS
%================================

\newenvironment{ctab}[2]{%
				\begin{spacing}{#1}
					\begin{center}
					\begin{tabular}{#2}}
					{\end{tabular}
					\end{center}
				\end{spacing}
				}


%================================
% Farbiger-Background-Box
%================================



\newcommand{\farbig}[1]{
\begin{tcolorbox}[%
				hbox,
				colback=miugray2!25,
				boxsep=-2pt,
				boxrule=0pt,
				arc=0pt,
				nobeforeafter,
				tcbox raise base,
				shrink tight,
				extrude by=1pt
				]
				{\color{miudiutex}{#1}}
\end{tcolorbox}~}

%================================
% SIMPLE SCHWARZER RAHMEN BOX
%================================


\newcommand{\boxfr}[1]{
	\begin{tcolorbox}[
	colframe=miugray2,
	colback=white!95!miugray2,
	boxrule = 0.5pt, 
%	toprule=0pt, bottomrule=0pt,rightrule=0pt,leftrule=0pt,
	boxsep=0pt,
	arc=4pt,
	drop shadow southeast,
	enhanced
	]
	\color{miudiutex}
	\begin{center}
	#1
	\end{center}
	\end{tcolorbox}
}

%================================
% Antwort-Box
%================================

\newcommand{\antwort}[1]{\par 
	\begin{tcolorbox}[
	colback=gray!15,
	breakable,
	enhanced,
	frame hidden,
	arc=5pt,
	boxrule = 2pt,
%	borderline west = {2pt}{0pt}{petunidin},
%	drop shadow=miugray2!30
	]
	\color{miudiutex}	#1
	\end{tcolorbox}
}

%================================
% Frame-IT
%================================

\newcommand{\frameit}[1]{\par 
	\begin{tcolorbox}[
	colback=gray!10,
	breakable,
	enhanced,
	frame hidden,
	arc=5pt,
	boxrule = 0pt,
	borderline west = {4pt}{0pt}{petunidin},
%	drop shadow=miugray2!30
	]
		#1
	\end{tcolorbox}
}



%================================
% Question BOX
%================================

\makeatletter
\newtcbtheorem{qstion}{Frage}{enhanced,
	breakable,
	colback=white,
	colframe=delphinidin!100,
%	capture=hbox,
	boxsep=-3pt,
	attach boxed title to top left={yshift*=-\tcboxedtitleheight},
	fonttitle=\bfseries,
%	title={#2},
	boxed title size=title,
	boxed title style={%
			sharp corners,
			rounded corners=northwest,
			colback=tcbcolframe,
			boxrule=0pt,
		},
	underlay boxed title={%
			\path[fill=tcbcolframe] (title.south west)--(title.south east)
			to[out=0, in=180] ([xshift=5mm]title.east)--
			(title.center-|frame.east)
			[rounded corners=\kvtcb@arc] |-
			(frame.north) -| cycle;
		},
	#1
}{def}
\makeatother



%================================
% SIMPLE Question BOX
%================================

\newcommand{\qssmall}{
\begin{tcolorbox}[hbox,colback=miugray2!55,boxsep=-4pt,boxrule=0pt,arc=2pt,nobeforeafter, tcbox raise base,shrink tight,extrude by=3pt]
\textbf{\color{white}{?}}
\end{tcolorbox}~~}

\newcommand{\folie}[1]{
\begin{tcolorbox}[%
				hbox,
				colback=miugray2!20,
				boxsep=2pt,
				boxrule=0.5pt,
				arc=2pt,
				nobeforeafter,
				tcbox raise base,
				shrink tight,
				extrude by=3pt]
				{\color{miudiutex}{Folie #1}}
\end{tcolorbox}~}


%================================
% NOTIZ BOX
%================================

\makeatletter
\newtcbtheorem{notizz}{Notiz}{enhanced,
	breakable,
	coltitle=white,
	colback=petunidin!10,
	colframe=petunidin!50,
	attach boxed title to top left={yshift*=-\tcboxedtitleheight},
	fonttitle=\bfseries,
%	title={#2},
	boxed title size=title,
	boxed title style={%
			sharp corners,
			rounded corners=northwest,
			colback=tcbcolframe,
			boxrule=0pt,
		},
	underlay boxed title={%
			\path[fill=tcbcolframe] (title.south west)--(title.south east)
			to[out=0, in=180] ([xshift=5mm]title.east)--
			(title.center-|frame.east)
			[rounded corners=\kvtcb@arc] |-
			(frame.north) -| cycle;
		},
		#1
}{def}
\makeatother



%================================
% Question BOX: Klausur
%================================

\makeatletter
\newtcbtheorem[no counter]{klausurinfo}{Klausurrelevant}{enhanced,
%	breakable,
	colback=miugray2!10,
	colframe=miugray2!75,
	attach boxed title to top left={yshift*=-\tcboxedtitleheight},
	fonttitle=\bfseries,
	title={#2},
	arc=4pt,
	boxrule=1pt,
	boxed title size=title,
	boxed title style={%
			sharp corners,
			rounded corners=northwest,
			colback=tcbcolframe,
			boxrule=0pt,
			arc=4pt,
		},
	underlay boxed title={%
			\path[fill=tcbcolframe] (title.south west)--(title.south east)
			to[out=0, in=180] ([xshift=5mm]title.east)--
			(title.center-|frame.east)
			[rounded corners=\kvtcb@arc] |-
			(frame.north) -| cycle;
		},
	#1
}{def}
\makeatother


%================================
% DEFINITION
%================================

\makeatletter
\newtcbtheorem[no counter]{defin}{Definition}{enhanced,
	breakable,
	colback=gray!10,
	colframe=gray,
	attach boxed title to top left={yshift*=-\tcboxedtitleheight},
	fonttitle=\bfseries,
	title={#2},
	arc=4pt,
	boxrule=1pt,
	fontlower=\footnotesize,
	collower=miugray2!50!miudiutex,
	boxed title size=title,
	boxed title style={%
			sharp corners,
			rounded corners=northwest,
			colback=tcbcolframe,
			boxrule=0pt,
			arc=4pt,
		},
	underlay boxed title={%
			\path[fill=tcbcolframe] (title.south west)--(title.south east)
			to[out=0, in=180] ([xshift=5mm]title.east)--
			(title.center-|frame.east)
			[rounded corners=\kvtcb@arc] |-
			(frame.north) -| cycle;
		},
	#1
}{def}
\makeatother

%================================
% Gesetzes-Text EU
%================================

\makeatletter
\newtcbtheorem[no counter]{gestex}{EU - Verordnung}{enhanced,
	breakable,
	colback=gray!10,
	colframe={gray!65!delphinidin},
	attach boxed title to top left={yshift*=-\tcboxedtitleheight},
	fonttitle=\bfseries,
	arc=4pt,
	boxsep=10pt,
	boxrule=1pt,
	fontlower=\footnotesize,
	collower=miugray2!50!miudiutex,
	boxed title size=title,
	overlay={%
			\node[xshift=-1cm,yshift=-1mm] at (frame.north east)
			{\includegraphics[width=8mm]{F://LaTeX-Files/shiet/eu}};
%			\node[draw,miugray2] at (frame.north) {{#2}}; 
			},%
	boxed title style={%
			sharp corners,
			rounded corners=northwest,
			colback=tcbcolframe,
			boxrule=0pt,
			arc=4pt,
		},
	underlay boxed title={%
			\path[fill=tcbcolframe] (title.south west)--(title.south east)
			to[out=0, in=180] ([xshift=5mm]title.east)--
			(title.center-|frame.east)
			[rounded corners=\kvtcb@arc] |-
			(frame.north) -| cycle;
		},
	#1
}{def}
\makeatother


%================================
% Gesetzes-Text Ger
%================================

\makeatletter
\newtcbtheorem[no counter]{gesgex}{Gesetzestext}{enhanced,
	breakable,
	colback=gray!10,
	colframe={gray!95!pelargonidin},
	attach boxed title to top left={yshift*=-\tcboxedtitleheight},
	fonttitle=\bfseries,
	title={#2},
	arc=4pt,
	boxrule=1pt,
	fontlower=\footnotesize,
	collower=miugray2!50!miudiutex,
	boxed title size=title,
	overlay={\node[xshift=-1cm,yshift=-3mm] at (frame.north east)
			{\includegraphics[width=8mm]{F://LaTeX-Files/shiet/ger}}; 
			},
	boxed title style={%
			sharp corners,
			rounded corners=northwest,
			colback=tcbcolframe,
			boxrule=0pt,
			arc=4pt,
		},
	underlay boxed title={%
			\path[fill=tcbcolframe] (title.south west)--(title.south east)
			to[out=0, in=180] ([xshift=5mm]title.east)--
			(title.center-|frame.east)
			[rounded corners=\kvtcb@arc] |-
			(frame.north) -| cycle;
		},
	#1
}{def}
\makeatother


%================================
% Question BOX: Wichtig
%================================

%\makeatletter
%\newtcbtheorem[no counter]{wichtig}{Wichtige Info}{enhanced,
%	breakable,
%	colback=pelargonidin2!80,
%	colframe=red2!100,
%	fonttitle=\bfseries,
%	title={#2},
%	leftrule=5mm,
%	boxed title size=title,
%	boxed title style={%
%			sharp corners,
%			rounded corners=northwest,
%			colback=tcbcolframe,
%			boxrule=0pt,
%		},
%	underlay boxed title={%
%			\path[fill=tcbcolframe] (title.south west)--(title.south east)
%			to[out=0, in=180] ([yshift=1mm,xshift=7mm]title.east)--
%			([yshift=1mm]title.center-|frame.east)
%			[rounded corners=\kvtcb@arc] |-
%			(frame.north) -| cycle;
%		},
%		watermark tikz={\draw[line width=2mm] circle (1.1cm)
%		node{\fontfamily{ptm}\fontseries{b}\fontsize{20mm}{20mm}\selectfont !};}],
%	attach boxed title to top left={xshift=5mm,yshift*=-\tcboxedtitleheight},
%	watermark zoom=0.65,
%	  underlay={%
%    \path[draw=none] (interior.south west) rectangle node[white]{\Huge \textbf{!}} ([xshift=-4mm]interior.north west);
%    },
%	#1
%}{def}
%\makeatother


\makeatletter
\newtcbtheorem[no counter]{wichtig}{Wichtige Info}{enhanced,
	breakable,
	colback=pelargonidin2!30,
	colframe=petunidin!100,
	fonttitle=\footnotesize\bfseries,
	title={#2},
	leftrule=4mm,
	boxed title size=title,
	boxed title style={%
			sharp corners,
			rounded corners=northeast,
			colback=tcbcolframe,
			boxrule=0pt,
		},
	underlay boxed title={%
			\path[fill=tcbcolframe] (title.south east)--(title.south west)
			to[out=180, in=0] ([yshift=2mm,xshift=-7mm]title.west)--
			([yshift=2mm]title.center-|frame.west)
			[rounded corners=\kvtcb@arc] |-
			(frame.north) -| cycle;
		},
		watermark tikz={\draw[line width=2mm] circle (1.1cm)
		node{\fontfamily{ptm}\fontseries{b}\fontsize{20mm}{20mm}\selectfont !};},
	attach boxed title to top right={yshift*=-\tcboxedtitleheight},
	watermark zoom=1.45,
	watermark opacity=0.35,
	  underlay={%
    \path[draw=none] (interior.south west) rectangle node[white]{\LARGE \textbf{!}} ([xshift=-4.3mm]interior.north west);
    },
    #1
}{def}
\makeatother





%================================
% ZITAT BOX
%================================

\tcbuselibrary{theorems,skins,hooks}
\newtcbtheorem{zitat}{Zitat}
{%
	enhanced,
	breakable,
	colback = gray!10!white,
	frame hidden,
	boxrule = 0sp,
	borderline west = {2pt}{0pt}{petunidin},
	sharp corners,
	detach title,
	before upper = \tcbtitle\par\smallskip,
	coltitle = gray!50!black,
	fonttitle = \bfseries\sffamily,
	description font = \mdseries,
	separator sign none,
	segmentation style={solid, gray!50!black},
}
{th}



\input{F://LaTeX-Files/shiet/vl.tex}
\usepackage{pgfplots}
\usepackage{booktabs, makecell}
\usepackage{awesomebox}
\usepackage{caption}
\usepackage{hypcap}
\usepackage[makeroom]{cancel}
%\setcounter{section}{1} % <--- FALLS KEINE SECTION DEFINIERT WIRD
    \pgfplotsset{
    	compat=1.18,
        every axis post/.style={
	yticklabel=\empty,
	xticklabel=\empty,
	ticks=none,
%	width=6cm,
	xlabel near ticks,
        },
    }
\usepgfplotslibrary{fillbetween}
\usetikzlibrary{patterns.meta}
\usetikzlibrary{calc,shapes.symbols,decorations.pathreplacing}




\begin{document}

\sidebar
\miushead{Vorlesung 09 + 10 - PLMT - KW 19 }{Vorlesung 09 + 10}
\vspace*{-4mm}
%\section{Vorlesung 09 - Wärme}
\subsection{Längenspezifischer Wärmestrom}



\begin{center}
\begin{tikzpicture}[scale=1]
\node(1)at(0,0){%
\begin{tikzpicture}
\draw[fill=miugray2!15,  postaction={pattern color=miugray2!50,pattern={Dots[angle=-45, distance=2mm,radius=0.7pt]}}] circle (2cm);
\draw[fill=white] circle (1cm);
\draw[->] (0,0) to (0:2.4cm)node[right] {$\dot{Q}$ = const.}; 
\draw[-,dashed] (0,0) to node[above] {$r_i$}(20:1cm) to node[above] {$r_a$}(20:2cm);
\node[yshift=-5mm](t1)at(0,0){$T_i$};
\node[yshift=-5mm](t1)at(0,-2){$T_a$};
\node[draw=none](t1)at(2,-2){$T_i > T_a$};
\end{tikzpicture}};
\node(2)at(8,0){
\begin{tikzpicture}
\draw[->](0,0) to (0,4) node[right] {$T$};
\draw[->](0,0) to (4,0) node[right] {$r$};

\draw[-] (0,3) node[left]{$T_i$}  to (1,3) to [bend right=35] (2,2) to (3.5,2) node[right] {$T_a$};
\draw[-,dashed]  (1,0) node[below] {$r_i$} to (1,4);
\draw[-,dashed]  (2,0) node[below] {$r_a$} to (2,4);
\node[align=left](3)at(4,3){logarithmischer Verlauf \\ der Temperatur};
\end{tikzpicture}
};

\node(3)at(4,0){
\begin{tikzpicture}[rotate=90,xscale=0.5,yscale=0.8]
\draw (0,0)-- (8,0) .. controls (8.25,0.15)and(8.25,0.85) .. (8,1) -- (0,1) [bend right=35]  to cycle;
\draw[<->] (0,-0.25) -- node[right] {$l$}(8,-0.25);
\draw (8,0) to [bend left=35] (8,1); 
\fill[miugray2!25] (8,0) [bend left=35] to (8,1) [bend left=35] to (8,0);
\draw[miugray2, fill=white] (8,0.5) ellipse (0.075cm and 0.25cm);
\end{tikzpicture}
};
\end{tikzpicture}
\captionof{figure}{Stationärer Fall der Wärmeleitung durch ein Rohr}
\end{center}
\begin{minipage}[t]{.4\textwidth}
\begin{tcbraster}[raster columns=1, raster valign=top]
\boxfr{
{\color{petunitext}\textbf{Herleitung:}}

\begin{align*}
d T &= - \frac{\dot{Q}}{\lambda \cdot A} \cdot d r \\
\int_{T_i}^T dT &= - \frac{\dot{Q}}{\lambda \cdot 2\pi l} \cdot \int_{r_i}^r \frac{dr}{r} \\
\left[T \right]_{T_i}^T &= - \frac{\dot{Q}}{\lambda \cdot 2\pi l} \cdot \left[ \ln(r) \right] _{r_i}^r \\
T-T_i &= - \frac{\dot{Q}}{\lambda \cdot 2\pi l} \cdot \left[ \ln(r) - \ln(r_i) \right] \\
	&= - \frac{\dot{Q}}{\lambda \cdot 2\pi l} \cdot \ln(\frac{r}{r_i}) \\
T_{(r)} &= T_i - \frac{\dot{Q}}{\lambda \cdot 2\pi l} \cdot \ln(\frac{r}{r_i}) \\
T_{(a)} &= T_i - \frac{\dot{Q}}{\lambda \cdot 2\pi l} \cdot \ln(\frac{r_a}{r_i}) \\
\dot{Q} &= \frac{2\pi l \cdot \lambda \cdot (T_i - T_a)}{\ln(\frac{r_a}{r_i})} \\ \\
\frac{\dot{Q}}{l} &= \dot{q}_L = (T_i - T_a) \cdot \frac{2 \pi \cdot \lambda}{\ln(\frac{r_a}{r_i})}
\end{align*}
}
\end{tcbraster}
\end{minipage}
\hfill
\begin{minipage}[t]{.5\textwidth}
\pp
{\color{petunitext}\textbf{Anwendung:}}
Berechnung des Wärmestroms für verschiedene Geometrien mittels der mittleren Fläche A

\begin{align*}
\dot{Q} = - \lambda \cdot \overline{A} \cdot \frac{dT}{dr}
\end{align*}
Für ein Rohr gilt:
\begin{align*}
\overline{A} = \frac{A_a-A_i}{\ln(\frac{A_a}{A_i})} = \frac{2 \pi \cdot l \cdot (r_a-r_i)}{\ln(\frac{r_a}{r_i})}
\end{align*}
\vspace*{-8mm}
\subsection{Konvektion}
{\color{petunitext}\textbf{Newtonsche Wärmeübergangsgleichung:}}
\begin{align*}
\dot{Q} = - \alpha \cdot A \cdot (T_M - T_W)
\end{align*}
$\dot{Q}$ = Wärmestrom durch Konvektion \par 
$T_W$ = Temperatur auf Wandoberfläche \par 
$T_M$ = Temperatur des Mediums in Strömung \par 
$A$ = Wandoberfläche \par 
$\alpha$ = Wärmeubergangskoeffizient
\end{minipage}


\pp
\reversemarginpar
\marginnote[\scriptsize \textbf{Turbulente Strömung} \par
Oberhalb einer sog. kritischen Geschwindigkeit geht die laminare in eine turbulente Strömung über, wobei Wirbelbewegungen und somit Kräfte entstehen, die auch entgegen der Bewegungsrichtung der Strömung wirken. Turbulente Strömung ist ein wesentlicher Faktor zur Erreichung einer guten Wärmeableitung durch Konvektion. Bei der Wärmeableitung über einen Kühlkörper kommt der Konvektion im Gegensatz zur Wärmestrahlung eine wesentlich größere Bedeutung zu. ]{}[-2cm]

\begin{center}
{\color{petunitext}\textbf{Konvektion von laminarer gegenüber turbulenter Strömung}}
\begin{tikzpicture}
\node(a)at(0,0){
\begin{tikzpicture}

\draw[fill=miugray2!15,  postaction={pattern color=miugray2!50,pattern={Dots[angle=-45, distance=2mm,radius=0.7pt]}}] (-0.2,0) rectangle (0.3cm,2cm);
\draw[->,miugray2] (0.6,0) to (0.6,0.3);
\draw[->,miugray2] (1,0) to (1,0.7);
\draw[->,miugray2] (1.4,0) to (1.4,1);
\draw[->,miugray2] (1.8,0) to (1.8,1.2);
\draw[-,dashed,miugray2] (0.3,0) to [bend left=25] (1.8,1.2);
\node[miugray2,align=left,anchor=west,scale=0.8](1)at(2,0.5){Geschwindigkeitsprofil \\ laminare Strömung};
\draw[petunitext,-] (0.3,0.5) node[left,xshift=1mm] {$T_1$} to (1.8,1.6) node[right] {$T_2$};
\draw[->,delphinidin] (1.8,1.8) to node [above] {$\dot{Q}$} (-0.5,1.8)  ;
\draw[->,black!60] (1.8,-0.1) to node [below] {x} (-0.5,-0.1)  ;
\end{tikzpicture}
};
\node(b)at(6,0){
\begin{tikzpicture}

\draw[fill=miugray2!15,  postaction={pattern color=miugray2!50,pattern={Dots[angle=-45, distance=2mm,radius=0.7pt]}}] (-0.2,0) rectangle (0.3cm,2cm);
\draw[->,miugray2] (0.6,0) to (0.6,0.3);
\draw[->,miugray2] (1,0) to (1,0.4);
\draw[->,miugray2] (1.4,0) to (1.4,0.4);
\draw[->,miugray2] (1.8,0) to (1.8,0.3);
\draw[-,dashed,miugray2] (0.3,0) to [bend left=35] (1.8,0.3);
\node[miugray2,align=left,anchor=west,scale=0.6](1)at(2,0.3){Geschwindigkeitsprofil \\ turbulente Strömung};
%\draw[brown] (0.3,0.5) to [bend left=30] (1.8,1.2) node[right,align=left,scale=0.7,xshift=2mm,yshift=-1mm] {hydrodynamische \\ Grenzschicht};
\draw[<->,brown] (0.3,0.5) to (0.8,0.5) node[left,xshift=1mm,yshift=3mm,align=left,anchor=west,scale=0.5] {hydrodynamische \\ Grenzschicht};
\draw[petunitext,-] (0.3,0.5) node[left,xshift=1mm] {$T_1$} to [bend left=30] (1.8,1.6) node[right] {$T_2$};
\draw[->,delphinidin] (1.8,1.8) to node [above] {$\dot{Q}$} (-0.5,1.8)  ;
\draw[->,black!60] (1.8,-0.1) to node [below] {x} (-0.5,-0.1)  ;
\draw[<->,black!60] (0.3,1.1) to (1.15,1.1) node[left,xshift=1mm,yshift=1.5mm,align=left,anchor=west,scale=0.5] {thermische \\ Grenzschicht};;
\end{tikzpicture}
};
\end{tikzpicture}
\par 
\textbf{Laminare Strömung:} Keine Konvektion in x-Richtung sondern Wärmeleitung
\end{center}

\pagebreak
\subsection{Wärmedurchgang durch eine Wand}

\begin{minipage}{.45\textwidth}
\begin{tikzpicture}[scale=0.8]
\draw[miugray2!75,fill=miugray2!25,dashed] (3.5,-0.5) rectangle (4.5,5.5); 
\draw[<->,black!65] (3.5,1) to node[above] {$d_W$} (4.5,1);
\node[miugray2!75,scale=1.5](a1)at(1.5,1.5){$\alpha_{Warm}$};
\node[miugray2!75,scale=1.5](a2)at(6.5,3.5){$\alpha_{Kalt}$};
\draw[->] (0,0) -- (0,5) node[right] {$T$}; 
\draw[->] (0,0) -- (8,0) node[right] {$x$}; 
\draw[->,line width=2pt,petunitext!65,rounded corners=3pt] (0,4) node[above,xshift=4mm] {$T_2$} .. controls (2.5,4) and (3,3.75) .. (3.5,3) to (4.5,2.5) .. controls (5,1.75) and (5.5,1.5) .. (8,1.5) node[right] {$T_2$};
\draw[-,dotted] (2.75,0) to (2.75,5);
\draw[-,dotted] (5.25,0) to node[right,rotate=90,yshift=2mm,scale=0.6,miugray2] {Thermische Grenzschicht} (5.25,5);
\end{tikzpicture}
%\captionof{figure}{Wärmedurchgang im stationären Fall}
\end{minipage}
\hfill
\begin{minipage}[t]{.45\textwidth}
\vspace*{-15mm}
\begin{align*}
\dot{Q} &= k \cdot A \cdot \Delta T_{ges}
\end{align*}
Wobei:
\begin{align*}
k &= (\frac{1}{\alpha_{Warm}} + \frac{d_W}{\lambda_W} + \frac{1}{\alpha_{Kalt}} )^{-1}
\end{align*}
\end{minipage}

\subsection{Konvektion: Prandtl-Zahl}
\begin{align*}
\text{Prandtl-Zahl} = \frac{\text{Impulstransport durch Reibung}}{\text{Wärmetransport durch Leitung}}
\end{align*}

\begin{tikzpicture}
\node(a)at(0,0){
\begin{tikzpicture}
\draw[fill=miugray2!15,  postaction={pattern color=miugray2!50,pattern={Dots[angle=-45, distance=2mm,radius=0.7pt]}}] (0,0) rectangle (-0.5,4); 
\draw[petunitext](0,1) .. controls (0.1,2.8) and (0.3,3)  .. (4,3);
\draw[<->,petunitext,dashed] (0,3.5) to (1,3.5);
\draw[<->,miugray2,dashed] (0,3.75) to (3.5,3.75);
\draw[miugray2](0,0.25) .. controls (0.25,0.75) and (2,1.7) .. (4,2);
\node[scale=1.1,font=\bfseries,yshift=3mm] at (3.5,0) {Pr $>$ 1};
\end{tikzpicture}
};
\node(b)at(6,0){
\begin{tikzpicture}
\draw[fill=miugray2!15,  postaction={pattern color=miugray2!50,pattern={Dots[angle=-45, distance=2mm,radius=0.7pt]}}] (0,0) rectangle (-0.5,4); 
\draw[miugray2,yshift=-1cm](0,1) .. controls (0.1,2.8) and (0.3,3) .. (4,3);
\draw[<->,miugray2,dashed] (0,3.75) to (1,3.75);
\draw[<->,petunitext,dashed] (0,3.5) to (3.5,3.5);
\draw[petunitext,yshift=1.5cm](0,0.25) .. controls (0.25,0.75) and (2,1.7) .. (4,2);
\node[scale=1.1,font=\bfseries,yshift=3mm] at (3.5,0) {Pr $<$ 1};
\end{tikzpicture}
};
\node[anchor=west,petunitext,scale=0.8](c)at(8.5,1){$\blacksquare$ Thermische Grenzschicht};
\node[anchor=west,miugray2,scale=0.8](d)at(8.5,0.5){$\blacksquare$ Hydrodynam. Grenzschicht};
\end{tikzpicture}
\end{document}
